%% kapa: exam=1 round=20 year=2009 period=1 form=A subject=economics subject_ko=경제원론 seq=13 number=53
%% source: Q-Net 감정평가사 기출문제 (cst00309) artlSeq=5024178 file=extracted/1008/5024178/2009년 제20회 감정평가사 1차시험 기출문제/1교시.hwp
%% conversion: kordoc 4.24.51 -> tools/gen1.py (auto); review=proofread-sonnet; flags=none
\begin{question}{53}
토마토케첩과 핫도그는 정상재이며, 서로 보완재이다. 핫도그 원료인 밀가루 가격 인상에 따른 핫도그 가격 상승의 효과로 옳은 것은?
\begin{choices}
  \item 토마토케첩의 균형가격 상승
  \item 핫도그의 균형공급량 증가
  \item 토마토케첩의 균형공급량 증가
  \item 핫도그의 균형수요량 증가
  \item 토마토케첩의 수요 감소
\end{choices}
\end{question}
\answer{5}
